Em 1989, a já extremamente popular linguagem C alcançou níveis de popularidade 
nunca antes vistos, sendo oficialmente reconhecida pelas organizações de 
padronização ANSI (americana) e ISO (internacional). Esse processo culminou na 
padronização conhecida como \texttt{C89} (ou ANSI C), que é o dialeto da 
linguagem que mais reconhecemos atualmente.

Atualizações e melhorias na padronização da linguagem são lançadas, geralmente, 
a cada 10 anos. Atualmente, as versões mais conhecidas e utilizadas no mercado 
são: \texttt{C89}, \texttt{C99}, \texttt{C11} e \texttt{C18}, sendo as versões 
\texttt{C99} e \texttt{C11} as mais famosas entre elas.

Apesar disso, há uma versão publicada em 31 de outubro de 2024 que trouxe 
algumas funcionalidades bem úteis e simples de usar. Gostaria de mencioná-las 
neste livro, pois de certa forma procuro trazer uma visão moderna da linguagem 
para iniciantes.

\subsubsection{Como ativar o padrão C23}

Antes de mostrar qualquer funcionalidade do padrão \texttt{C23} (também chamado 
de \texttt{C2X}), é necessário explicar como ativá-lo no compilador. 
Diferentemente de linguagens modernas como JavaScript ou Python, onde 
funcionalidades são adicionadas por meio de pacotes, podemos ativar as 
funcionalidades do \texttt{C23} usando uma \texttt{flag} de compilação que 
indica ao GCC/Clang qual padrão desejamos, da seguinte forma:

\begin{textbox}
    gcc -o main main.c -std=c2x
    \\ ou
    gcc -o main main.c -std=c23
\end{textbox}

O comando pode variar um pouco dependendo do compilador, mas no geral deve ser 
como mostrado acima. A opção \texttt{-std=} também pode ser utilizada para 
especificar qualquer outro padrão da linguagem. Portanto, seria possível usar 
comandos como:

\begin{textbox}
    gcc -o main main.c -std=c18
    gcc -o main main.c -std=c11
    gcc -o main main.c -std=c99
    gcc -o main main.c -std=c89
\end{textbox}

Qualquer um desses é válido como especificação do padrão a ser utilizado pelo 
compilador. Isso é bastante útil em projetos mais antigos ou bem consolidados em 
um padrão específico da linguagem, sendo necessário, nesses casos, indicar ao 
compilador qual padrão utilizar.

\subsubsection{Separador de dígitos}

O separador de dígitos é apenas um recurso visual para melhorar a legibilidade 
de números muito grandes, que dificilmente seriam lidos sem nenhum tipo de 
separador. Trata-se de uma função meramente visual, sem qualquer implicação 
lógica no código.

\begin{excode}
#include <limits.h>
#include <stdio.h>

int main(void)
{
    int x = 19'251'916; // mesmo que 19251
    long int y = 36'854'775'807'192; // mesmo que 36854775807

    printf("int: %i\\n", x);
    printf("long int: %li\\n", y);

    return 0;
}
\end{excode}

\subsubsection{Literais binários}

Em algumas aplicações de baixo nível, por vezes é necessário declarar valores 
explicitamente bit a bit. Anteriormente, isso precisava ser feito usando 
hexadecimal ou octal, o que nem sempre era intuitivo.

\begin{excode}
...
unsigned char mascara = 0b11000101;
unsigned char verdadeiro = 0b1;
unsigned char falso = 0b0;
...
\end{excode}

Agora, utilizando o prefixo \texttt{0b} ou \texttt{0B} é possível definir 
valores em binário explicitamente, assim como mostrado acima.

\subsubsection{Ponteiro nulo}

Esta é outra mudança que traz melhorias de legibilidade, porém com um aspecto 
um pouco mais técnico. Desde a versão C89 da linguagem, a macro \texttt{NULL} 
tem sido definida como:

\begin{excode}
...
\#define NULL ((void*)0)
...
\end{excode}

No geral, essa definição é satisfatória para o valor nulo. Porém, em algumas 
plataformas essa mesma macro acaba sendo definida simplesmente como \texttt{0}, 
o que gera problemas de portabilidade. Um código que depende fortemente da 
definição de \texttt{NULL} pode quebrar em uma plataforma onde ela é definida 
de maneira diferente.

Um exemplo simples de como isso poderia acontecer é com a seguinte função:

\begin{excode}
...
printf("\%p\n", NULL);  // possui comportamento indefinido
...
\end{excode}

Caso \texttt{NULL} esteja definido como um inteiro de valor \texttt{0}, o 
comportamento é indefinido e pode resultar em lixo de memória, o que é algo 
certamente indesejado. Assim, a constante \texttt{nullptr} foi introduzida para 
corrigir esse problema:

\begin{excode}
...
int* x = nullptr;
...
\end{excode}

Agora temos um ponteiro para inteiro recebendo um ponteiro nulo genérico de 
maneira \textbf{explícita}. Dessa forma, o compilador consegue diferenciar 
\texttt{nullptr} de \texttt{((void*)0)} ou de um inteiro de valor \texttt{0}. 
Isso resolve o problema de forma explícita, forçando o programador a deixar 
claro o que está sendo feito no código e também resolvendo questões de 
compatibilidade com outros sistemas e de interoperabilidade com C++.

\subsubsection{Formato de binário para saídas/entradas}

Em programação de sistemas de baixo nível, é comum ter que lidar com 
manipulação ou visualização de dados binários. Anteriormente, a linguagem C 
necessitava de funções implementadas manualmente para realizar tais tarefas. 
Agora, podemos usar o especificador \texttt{\%b} para exibir valores em binário 
de forma nativa:

\begin{excode}
#include <stdio.h>

int main(void) {
    int i = 67;
    printf("i: %b\n", i); // Exibe "i: 1000011"
    int j = 240;
    printf("j: %b\n", j); // Exibe "j: 11110000"
    int k = 101;
    printf("k: %b\n", k); // Exibe "k: 1100101"

    return 0;
}
\end{excode}

\subsubsection{True e False}

Apesar de ser uma funcionalidade \textbf{muito} simples presente na maioria das 
linguagens, C não suporta variáveis booleanas nativamente desde suas primeiras 
versões. Por conta disso, sempre que era necessário implementar um sistema de 
\texttt{flags} usando variáveis booleanas, era preciso incluir a biblioteca 
\texttt{<stdbool.h>} para usar as palavras-chave \texttt{true}, \texttt{false} 
e \texttt{bool}.

Agora, com esse recurso nativo da linguagem, é possível fazer:

\begin{excode}
#include <stdio.h>

int main(void) {
    bool v = true;
    bool f = false;

    if (v == true)
        printf("v: %i (em binário = %02b)\n", v, v);
    if (f == false)\
        printf("f: %i (em binário = %02b)\n", f, f);

    return 0;
}
\end{excode}

Tudo isso sem se preocupar com inconsistências ao incluir uma biblioteca para 
algo tão fundamental.

\subsubsection{\_BitInt(N)}

Em algumas aplicações de baixo nível como protocolos de rede, \texttt{hardware}, 
processamento de áudio, criptografia etc. É necessário trabalhar com tipos 
inteiros cuja largura não é potência de dois (por exemplo, 7 bits, 13 bits). 
Simular esses tipos com \texttt{structs} e máscaras era trabalhoso. Usando o 
novo tipo \texttt{\_BitInt(N)} podemos fazer:

\begin{excode}
#include <stdio.h>

int main(void) {
    \_BitInt(5) i = 0;
    i = 15;
    printf("i: %08b\n", i);
    i = -16;
    printf("i: %08b\n", i);

    return 0;
}
\end{excode}

Neste exemplo, declaramos um inteiro de 5 bits com sinal. O bit mais 
significativo é o bit de sinal, e os 4 bits restantes representam a magnitude 
(em complemento para dois). Assim, o intervalo de valores representáveis vai de 
-16 a 15.

\subsubsection{typeof e typeof\_unqual}

Ambos funcionam de forma bem semelhante, mas com uma diferença sutil que possui 
aplicações práticas na implementação de macros. A palavra-chave \texttt{typeof} 
é um operador que obtém o tipo de uma expressão em tempo de compilação. Isso 
significa que é possível fazer algo como:

\begin{excode}
#include <stdio.h>

int main(void) {
    int i = 10;
    typeof(i) j = i; // Mesmo que "int j = i;"

    printf("i: %i\n", i); // Exibe "i: 10"
    printf("j: %i\n", j); // Exibe "j: 10"

    return 0;
}
\end{excode}

O operador \texttt{typeof} sinaliza ao compilador que ele deve substituir 
\texttt{typeof(i)} pelo tipo da variável passada como parâmetro.

Porém, em alguns casos onde se deseja criar variáveis do mesmo tipo que outra, 
pode ser indesejado trazer qualificadores como \texttt{const}, \texttt{volatile} 
ou \texttt{restrict}. É aí que \texttt{typeof\_unqual} se mostra útil. Com ele 
podemos fazer:

\begin{excode}
#include <stdio.h>

int main() {
    const int a = 10;
    volatile double b = 3.14;

    // typeof mantém os qualificadores
    typeof(a) x = a;        // x é const int
    typeof(b) y = b;        // y é volatile double

    // typeof_unqual remove os qualificadores
    typeof_unqual(a) u = a; // u é int (sem const)
    typeof_unqual(b) v = b; // v é double (sem volatile)

    // Possível modificar u e v sem problemas
    u = 20;
    v = 2.71;

    printf("u = %d, v = %f\n", u, v);
    return 0;
}
\end{excode}

\begin{postscript}
Com esses operadores é possível criar macros como:

\begin{excode}
...
#define SWAP(a,b) { typeof(a) tmp = a; a=b; b=tmp; }
...
\end{excode}

Essa macro recebe dois valores, utiliza \texttt{typeof} para criar uma variável 
\texttt{tmp} e troca os valores entre \texttt{a} e \texttt{b}. O uso de 
\texttt{typeof\_unqual} é desnecessário nesse caso, pois não é desejado tentar 
modificar variáveis com qualificadores \texttt{const} ou \texttt{restrict}.
\end{postscript}

\subsubsection{deprecated e nodiscard}

Os atributos \texttt{deprecated} e \texttt{nodiscard} têm usos focados 
em legibilidade e documentação, sendo importantes principalmente para explicitar 
o comportamento de funções e tipos. Vamos entender cada um.

\texttt{deprecated} sinaliza que um elemento (função, variável, 
\texttt{struct}, etc.) está obsoleto e não deve ser usado em novos código, 
embora ainda funcione. Pode incluir uma mensagem opcional.

\begin{excode}
#include <stdio.h>

[[deprecated("Use a função 'soma_nova()' em vez disso")]]
int soma_antiga(int a, int b)
{
    return a + b;
}

[[deprecated]]
int soma_velhaca(int a, int b)
{
    return a + b;
}

...

int main(void) {
    int a = 10;
    int b = 15;

    printf("soma: %i\n", soma_antiga(a, b));
    printf("soma: %i\n", soma_velhaca(a, b));

    return 0;
}
\end{excode}

No código acima, temos duas funções obsoletas. Em bases de código antigas, é 
comum sinalizar funções que ainda funcionam mas serão removidas futuramente por 
meio de comentários. O atributo \texttt{deprecated} torna essa sinalização 
explícita e gera \texttt{warnings} durante a compilação, incentivando a migração 
para alternativas.

Observe que o atributo pode ser usado com ou sem mensagem entre aspas. O 
resultado da compilação do código acima exibe o seguinte:

\begin{textbox}
main.c: In function 'main\':                                                                                                                                                                       
main.c:19:5: warning: 'soma_antiga' is deprecated: Use a função
    'soma_nova()' em vez disso [-Wdeprecated-declarations]
   19 |     printf("soma: %i\n", soma_antiga(a, b));
      |     ^~~~~~
main.c:4:5: note: declared here
    4 | int soma_antiga(int a, int b)
      |     ^~~~~~~~~~~
main.c:20:5: warning: 'soma_velhaca' is deprecated
    [-Wdeprecated-declarations]
   20 |     printf("soma: %i\n", soma_velhaca(a, b));
      |     ^~~~~~
main.c:10:5: note: declared here                                                                                                                                                                  
   10 | int soma_velhaca(int a, int b)
      |     ^~~~~~~~~~~~
\end{textbox}

Já o atributo \texttt{nodiscard} indica que o valor retornado por uma 
função não deve ser ignorado. Se o chamador descartar o valor, o compilador 
emitirá um aviso. Isso é útil para funções que alocam recursos, retornam códigos 
de erro ou produzem resultados essenciais.

\begin{excode}
#include <stdio.h>

[[nodiscard]]
int soma_nova(int a, int b)
{
    return a + b;
}

int main(void) {
    int a = 10;
    int b = 15;

    soma_nova(a, b); // Valor não é atribuído a nenhuma variável

    return 0;
}
\end{excode}

No caso acima, a função \texttt{soma\_nova()} é chamada, mas o valor retornada 
pela mesma não é atribuída a nenhuma ou usado de nenhuma forma, o que faz com 
que o compilador retorne um erro como:

\begin{textbox}
main.c: In function 'main':
main.c:12:5: warning: ignoring return value of 'soma_nova', declared with 
    attribute 'nodiscard' [-Wunused-result]
   12 |     soma_nova(a, b); // Valor não é atribuído a nenhuma variável
      |     ^~~~~~~~~~~~~~~
main.c:3:19: note: declared here
    3 | [[nodiscard]] int soma_nova(int a, int b)
      |                   ^~~~~~~~~
\end{textbox}

Dito isso, é notável como ambos os atributos ajudam a tornar o código mais 
robusto e auto-documentado.

\begin{postscript}
É possível escrever mensagens de erro personalizadas da mesma forma que é feito 
com o atributo \texttt{deprecated}.

\begin{excode}
[[nodiscard("O valor de soma_nova() foi descartado")]]
\end{excode}

Desta forma, caso o valor retornado pela função não seja usado devidamente, a 
mensagem de erro personalizada deve aparecer como um erro durante o processo 
de compilação.

\end{postscript}

\subsubsection{\#elifdef e \#elifndef}

Essas duas diretivas foram introduzidas apenas para tornar a definição 
condicional de macros mais conveniente. Antes, para escrever código que 
dependesse do sistema operacional, por exemplo, era necessário algo como:

\begin{excode}
#include <stdio.h>

int main(void)
{
    #ifdef _WIN32
        printf("Executando no Windows\n");
    #elif defined(__linux__)
        printf("Executando no Linux\n");
    #elif defined(__APPLE__)
        printf("Executando no macOS\n");
    #else
        printf("Sistema operacional desconhecido\n");
    #endif

    return 0;
}
\end{excode}

Com as novas diretivas, o mesmo código pode ser reescrito de forma mais limpa e 
direta:

\begin{excode}
#include <stdio.h>

int main(void)
{
    #ifdef _WIN32
    printf("Executando no Windows\n");
    #elifdef __linux__
        printf("Executando no Linux\n");
    #elifdef __APPLE__
        printf("Executando no macOS\n");
    #else
        printf("Sistema operacional desconhecido\n");
    #endif

    return 0;
}
\end{excode}

A diretiva \texttt{\#elifdef} é equivalente a \texttt{\#elif defined(...)}, e 
\texttt{\#elifndef} equivale a \texttt{\#elif !defined(...)}. Elas tornam o código 
mais legível, especialmente em cadeias longas de condições.

\subsubsection{\#embed}

Esta é, provavelmente, uma das funcionalidades mais úteis, pelo menos na minha 
opinião, dentre todas as apresentadas. Trata-se de uma diretiva de 
pré-processamento que permite incluir dados binários diretamente no executável 
final gerado pela compilação.

Antes, era comum usar ferramentas como \texttt{xxd}, \textit{scripts} para gerar 
arrays em C ou até mesmo a função \texttt{fopen()} para ler arquivos durante a 
execução. A nova diretiva, no entanto, resolve tudo em tempo de compilação, sem 
custos de leitura em tempo de execução. Além disso, integra-se bem com sistemas 
de \textit{build}, que podem detectar alterações no arquivo incorporado e 
recompilar automaticamente. Tudo isso de forma bem simples, como no exemplo:

\begin{excode}
// Incluindo um arquivo de shaders que deve ser utilizado pelo programa
static const char arquivo_shaders[] = {
    #embed "shaders/vertex.glsl"
    , '\0'  // garante terminação nula para o array de bytes do arquivo
};
\end{excode}

O código acima incorpora o arquivo \texttt{"shaders/vertex.glsl"} no binário 
gerado pelo compilador e adiciona um caractere nulo ao final do array para 
garantir que a string seja terminada corretamente.

\subsubsection{memset\_explicit()}

Essa função foi adicionada para resolver problemas de segurança criados por uma 
interação, um tanto contraditória, entre o compilador e a função 
\texttt{memset()}. A função \texttt{memset()} é geralmente usada para limpar 
áreas de memória que armazenaram informações sensíveis. Por exemplo:

\begin{excode}
#include <string.h>

void processa_senha(const char *senha)
{
    char buffer[64];
    strcpy(buffer, senha);
    // ... usa a senha para algo (ex.: autenticar conta)

    // Tenta limpar a senha da memória
    memset(buffer, 0, sizeof(buffer));
}
\end{excode}

Nesse exemplo, após a chamada de \texttt{memset()}, a variável \texttt{buffer} 
não é mais usada. O compilador, ao otimizar o código com as flags \texttt{-O2} 
ou \texttt{-O3}, pode remover completamente a chamada de \texttt{memset()}, pois 
considera que escrever em uma variável que logo será descartada é redundante. 
Com isso, a senha permanece na pilha até ser sobrescrita por outra função, o que 
pode ser explorado por usuários mal-intencionados para recuperar dados sensíveis 
que ainda estão na memória.

A função \texttt{memset\_explicit()} resolve exatamente esse problema: ela foi 
projetada para ser imune a otimizações do compilador, ou seja, será chamada 
independentemente do nível de otimização usado. O resultado é um código mais 
seguro. Basta substituir \texttt{memset()} por \texttt{memset\_explicit()}:

\begin{excode}
#include <string.h>

void processa_senha(const char *senha)
{
    char buffer[64];
    strcpy(buffer, senha);
    // ... uso da senha

    // Garante a limpeza do buffer da senha
    memset_explicit(buffer, 0, sizeof(buffer));
}
\end{excode}

\begin{postscript}
No geral, estas são as adições mais relevantes trazidas pelo novo padrão C23. O 
foco dessa padronização foi remover funcionalidades obsoletas, melhorar a 
segurança, a portabilidade e a compatibilidade com C++.

Dentro das limitações da linguagem e do que fazia sentido ser modificado ou 
adicionado, é possível dizer que o novo padrão trouxe melhorias 
significativamente boas para o C.

No entanto, como já dito, aqui foram abordadas apenas as mudanças mais 
relevantes. Caso você queira saber \textbf{tudo} o que foi adicionado, recomendo 
fortemente procurar a lista completa de alterações e lê-la por conta própria. 
Como programador C, é importante ter a capacidade de buscar fontes de 
conhecimento e saber como se atualizar.
\end{postscript}
